\cotefile{156-6}
\scanpage{34}\mspage{33}

\uline{At$^{*}$2} Si $\Phi,\Psi\in{\unsure{\mathcal{F}|\mathscr{F}}}_M$, et $A\in\Phi\cap\Psi$, alors
$\Phi_A=\deleted{\unsure{\Phi}}\Psi_A$; i.e. la multistrate définie par $A$, soit dans
$\deleted{\unsure{\in\Phi}}\Phi$, soit dans $\Psi$, est \unsure{le|la} $=$.
% slanted
\margin{At$^{*}$1 et At$^{*}$2 impliquent \unsure{déjà} C\,1, C\,2 pour $\mathcal F$ $(\mathcal F_M,\leq\overset{\unsure{\mathrm{dfn}|\mathrm{déf}}}{=}\subset)$}

On l\deleted{'}a\deleted{ppellera} notera. $\deleted{\unsure{X|A}}\widetilde A$. ``Toute strate porte canoniquement une structure de multistrate'', indépendamment du \deleted{\unsure{le}} $\Phi\in\mathcal F_M$ dont c'est une strate.

\tikz[remember picture,overlay]\coordinate (msAtFour);
\uline{At$^{*}$\deleted{\unsure{3}}4} Si $\Phi,\Psi\in\mathcal F_M$ sont telles que $\forall\,A\in\Phi$, $B\in\Psi$, $\deleted{\unsure{\Phi_A\cup}}$ on ait $\Phi_A\cup\Psi_B\in\mathcal F_M$,
\margin{\uline{NB} il ne suffit \uline{pas} que $\Phi_A,\Psi_B$ soient compatibles \unsure{comme} figures ensemblistes!}
alors $\Phi\cup\deleted{\unsure{\Phi}}\Psi\in\mathcal F_M$.

\uline{At$^{*}$\deleted{\unsure{4}}5} L'ensemble $\deleted{\unsure{\mathcal{F}}}M^{*}$ de toutes les \deleted{\unsure{multi|mult}}strates \inserted{$A$} de \deleted{toutes} les $\Phi\in\mathcal F_M$ est une \unsure{figure} dans $M$ de support $M$. De plus,
$x,y\in M$, $x\ne y$, existe $A\in M^{*}$ telle que $x\in A$, $y\notin A$ ou $y\in A$, $x\notin A$.
($M^{*}$ ``sépare les points'' de $M$).
% slanted
\margin{i.e. tout $x\in M$ est contenu dans une strate \unsure{minimale|minim.}, les strates séparent les points, et \deleted{\unsure{pour}} \deleted{\unsure{une}} les \deleted{\unsure{strates}} strates \unsure{n'est|n'ont|ne sont} pas \unsure{réunion|réunions} des strates strictement \deleted{\unsure{plus}} contenues dans elles, (pour $\subset$)}
% The long arrow leads from the bracket at At*3 to the space just above At*4.

\tikz[remember picture,overlay]\coordinate (msAtThree);
\uline{At$^{*}$3} Soient $A,B$ deux strates relatives à $\mathcal F_M$ (i.e. des strates de $\Phi,\Psi\in\mathcal F_M$), telles que $A\subset B$, alors $\widetilde A\ll\widetilde B$\tikz[remember picture,overlay]\coordinate (msAtEnd);
\figure{\tikz[remember picture,overlay]{
  \draw ([xshift=7pt,yshift=8pt]msAtThree) .. controls +(-14pt,1pt) and +(-4pt,5pt) .. ([xshift=-5pt]msAtThree)
        .. controls +(-8pt,-7pt) and +(-13pt,0pt) .. ([xshift=7pt,yshift=-3pt]msAtThree |- msAtEnd);
  \draw[->] ([xshift=-5pt]msAtThree) .. controls +(-26pt,12pt) and +(-16pt,-12pt) .. ([yshift=8pt]msAtFour);
}}
% margin note written sideways
\margin{NB On \unsure{pourrait|p. maintenant} bloquer At$^{*}$2 et At$^{*}$3 en un seul : $\Phi,\Psi\in\mathcal F_M$, deux $A\in\Phi$, $B\in\Psi$, $A\subset B$ $\Longrightarrow\Phi_A\ll\Psi_B$. \uline{Cor.} Si $A=B$, $\Phi_A=\Psi_{\unsure{A|B}}$ d'où $A\mapsto\widetilde A$\,\ldots}

J'espère que je n'ai pas oublié d'axiomes!
